\section{Exact two-class ambiguity for all parameters}\label{sec:sharp-ambiguity}
\label{sec:rigidity}
Fix integers $a\ge0$, $b\ge1$, and $D=a+b\ge3$. The planted classes
$V_+,V_-$ have size $N$. Their internal factors are independent uniform
simple $a$-regular graphs, and the independent cross factor is uniform
simple $b$-regular bipartite. Degree-zero factors are absent. Limits
run through every admissible size, with $aN$ even and the simple factors
existing. Put $n=2N$, $\delta=a-b$, and
$\sigma=(\one_N,-\one_N)$. Let $K(G)$ count valid unordered
$(a,b)$ bisections, including the plant. Define
\begin{equation}\label{eq:ambiguity-exponent}
 \begin{split}
 m&=\begin{cases}\delta+1,&\delta\ge0,\\-\delta,&\delta<0,\end{cases}
 \qquad \nu=(D-2)m,\\
 C_{a,b}&=\begin{cases}
 \displaystyle\frac{((bt)!)^2(a!)^{2t}}
 {(t!)^2(b!)^{2t}a^{t\delta}},&\delta\ge0,\quad t=\delta+1,\\[6pt]
 \displaystyle\frac{((as)!)^2(b!)^{2s}}
 {(s!)^2(a!)^{2s}b^{s^2}},&\delta<0,\quad s=-\delta.
 \end{cases}
 \end{split}
\end{equation}
The first expression has $a\ge b\ge1$; at $\delta=0$ it has $a=b\ge2$.
The second is well defined at $a=0$ using $0!=1$.

\begin{theorem}[Boolean rigidity]\label{thm:boolean}
For every such fixed parameter pair,
\begin{equation}\label{eq:boolean-rate}
 \Prob\{\exists\tau\ne\pm\sigma:\tau\in\{-1,1\}^{2N},\quad
                         A\tau=\delta\tau\}=O_{a,b}(N^{-\nu}).
\end{equation}
\end{theorem}

\begin{theorem}[Exact ambiguity and Bayes error]\label{thm:sharp-ambiguity}
Along every admissible size sequence,
\begin{equation}\label{eq:sharp-ambiguity}
 \Prob\{K(G)\ge2\}\sim C_{a,b}N^{-\nu},\qquad
 \Prob\{K(G)\ge3\}=O_{a,b}(N^{-\nu-1}).
\end{equation}
Under a uniformly hidden balanced unordered plant, optimal Bayes error
is $R_N\sim(C_{a,b}/2)N^{-\nu}$. An always-correct Las Vegas complete
listing, followed by selection of any member, attains $R_N$ exactly
in unconditional expected $\widetilde O_{a,b}(n^2)$ bit work and
expected $\widetilde O_{a,b}(n)$ workspace. A deterministic estimator
with bounded $O_{a,b}(n^3+n^{\nu+2})$ elementary work and $O(n^2)$
storage attains the same leading Bayes constant.
\end{theorem}

These constants use block size $N$. For example, $(1,2)$ gives
$(\nu,C)=(1,2)$, $(1,3)$ gives $(4,16)$, $a=b\ge2$ gives
$(\nu,C)=(2a-2,(a!)^2)$, and $a=0,b\ge3$ gives
\[
 \nu=b(b-2),\qquad C_{0,b}=\frac{(b!)^{2b-2}}{b^{b^2}}.
\]
The assortative cubic pair $(2,1)$ has $(\nu,C)=(2,4)$.

\subsection{Minimum balanced flips and all larger supports}
Since $A\one=D\one$ and $D-\delta=2b>0$, every integer sign
eigenvector at $\delta$ is balanced. An alternative therefore flips
equally many $j$ vertices in each old class. Identifying global sign
lets us take $j\le N/2$. Write
\[
 w=(\sigma-\tau)/2=\one_{S_+}-\one_{S_-},\qquad
 |S_+|=|S_-|=j,\qquad Aw=\delta w.
\]
At a supported vertex, $\delta=i-c$, where $i$ is its number of
same-old-side support neighbors and $c$ its number of opposite-side
support neighbors. If $\delta\ge0$, simplicity gives
$\delta\le j-1$. Thus $j\ge t=\delta+1$, and equality forces two
$K_t$ internal cores and no support-cross edge. Every support vertex
then has $b$ internal and $b$ cross exterior edges. This includes
$\delta=0$, when both cores are singletons. If $\delta=-s<0$,
then $-s=i-c\ge-j$, so $j\ge s$. Equality forces no internal
support edge and the complete support core $K_{s,s}$; every support
vertex has $a$ internal and $a$ cross exterior edges.
These minimum assertions are used only in the simple graph law.

We bound all $j\ge m+1$ uniformly. At a zero coordinate of $w$,
exactly one incident edge from its nonzero support is impossible.
For a fixed support $S$ of size $h$, explore all incident pairings
in the unconditioned typed configuration model. A collision means
that the partner lies in $S$ or at a previously touched exterior
vertex. If $e$ edges lie in $S$, $B$ leave it, and $z$ exterior
vertices are touched, then
\[
 Dh=2e+B,\qquad z\le B/2,\qquad
 \#\text{collisions}=e+B-z\ge Dh/2.
\]
For $h\le\varepsilon N$, all present pools retain $\Theta_{a,b}(N)$
stubs and only $O_D(h)$ eligible collision stubs have appeared.
Each conditional collision probability is at most $C_{a,b}h/N$.
Absent internal pools are omitted, so no division by $a$ is used at
$a=0$. Degree-one matchings obey the same conditional exposure law.
Padding to $Dh$ steps with dummy noncollisions and union bounding
the collision positions gives
\[
 \Prob\{S\text{ has the necessary exterior condition}\}
 \le2^{Dh}(C_{a,b}h/N)^{Dh/2}.
\]
After summing over $\binom{2N}{h}$ supports this is at most
\begin{equation}\label{ap:eq:support}
 f_h=\left[L_{a,b}(h/N)^\eta\right]^h,
 \qquad \eta=(D-2)/2.
\end{equation}
Its consecutive ratio is at most
$L_{a,b}e^\eta((h+1)/N)^\eta$. Choose a small fixed
$\varepsilon$ making this at most one-half. The sum starting at
$h=2m+2$ is $O(N^{-(D-2)(m+1)})=O(N^{-\nu-1})$.
Conditioning the actual factors simple costs a bounded positive
factor, as justified below.

For macroscopic flips use real one-hot labels, not a sign-variance
argument requiring an internal slot. The candidate quotient is
$\left(\begin{smallmatrix}a&b\\b&a\end{smallmatrix}\right)$.
The exact count \eqref{eq:fmp-exact-types} applies to these integer
local equations. Symmetrization only within each old neighbor type
preserves root and pooled edge laws and increases the exponent.
Let $\Phi$ be the symmetrized exponent per physical vertex, with
old weights $\alpha_+=\alpha_-=1/2$, internal weights $a/4$, and
cross weight $b/2$. The regular frame has the symmetric omission
certificate $q_{ij}=1/D$ on every present factor. The oriented flip
fraction $r=j/N\le1/2$ is exactly the defect from the closest
old-block point map. Thus \eqref{eq:unknown-gap}, with
$c=1-2/D$, gives
\begin{equation}\label{ap:eq:macro-gap}
 2N\Phi\le-c_DNr^2,\qquad
 c_D=\frac{9(D-2)^3}{2D^7}>0.
\end{equation}
This gap is valid for $a=0,1$ and singular candidate quotients,
including $\delta=0$. Polynomially many fixed-alphabet types and
the uniform factorial remainders make every range $r\ge\varepsilon/2$
exponentially small. Combining it with \eqref{ap:eq:support} proves
\begin{equation}\label{ap:eq:large}
 \Prob\{\text{an alternative with least-oriented }j\ge m+1\}
                         =O_{a,b}(N^{-\nu-1}).
\end{equation}
The minimum $2m$ support size here comes from balanced flips; no
corresponding assertion about arbitrary real eigenvectors is used.

\subsection{All minimum patterns and their exact multiplicity}
Put $c=b$ in the nonnegative-$\delta$ case and $c=a$ otherwise.
At an exterior old-plus vertex, $Aw=\delta w$ requires equal
numbers of internal neighbors in $S_+$ and cross neighbors in $S_-$;
denote this common number by $m_v$. The same holds on the other
side, and on each side $\sum_v m_v=cm$. Hence each touched exterior
set has at most $cm$ vertices. A minimum flip is leading when both
have $cm$ vertices, so every positive $m_v=1$.

Its complete support-incident motif $H$ has
\[
 v_0=2m+2cm,\qquad e_0=2e_{\rm int}+e_{\rm cross},
 \qquad e_0-v_0=\nu,
\]
where, per internal factor,
\[
 \begin{array}{c|cc}
 &e_{\rm int}&e_{\rm cross}\\ \hline
 \delta\ge0&t\delta/2+bt&2bt\\
 \delta<0&as&s^2+2as=sD
 \end{array}
\]
A nonleading minimum pattern merges at least one exterior vertex
while preserving all $e_0$ distinct incident edges. There are finitely
many feasible simple typed patterns, each with at most $v_0-1$
vertices. Their placements and prescribed-stub probabilities give
$O(N^{-\nu-1})$, even after bounded-cost simplicity conditioning.
At $c=0$ the exterior is empty and all minimum patterns are leading.

Choose $S_+,S_-$ of size $m$ and disjoint exterior sets $W_+,W_-$
of size $cm$ on their old sides. Four maps specify the boundary:
\[
 W_+\to S_+\text{ internally},\quad W_+\to S_-\text{ across},\quad
 W_-\to S_-\text{ internally},\quad W_-\to S_+\text{ across}.
\]
Each has exactly $c$ preimages at every support vertex. Thus the
exact number of leading data choices is
\begin{equation}\label{ap:eq:data}
 A_N=\binom Nm^2\binom{N-m}{cm}^2
       \left[\frac{(cm)!}{(c!)^m}\right]^4
 \sim N^{v_0}\frac{((cm)!)^2}{(m!)^2(c!)^{4m}}.
\end{equation}
Empty maps count one. A realized leading flip determines its exterior
sets and all maps uniquely. Conversely these edges saturate every
support vertex and balance each exterior contribution, so they give
a valid flip. Other exterior edges do not matter. For $N>2m$ only
one orientation of an alternative has $m$ flips, so there is no
extra orientation or drawing-automorphism factor.

For $c>0$ the number of physical stub assignments for one data choice is
\[
 c_H=(a!b!)^{2m}(ab)^{2cm}.
\]
All support stubs are assigned to distinct prescribed neighbors, so
these assignments are disjoint events. One has exact configuration
probability
\begin{equation}\label{ap:eq:prescribed}
 P_N^*=\prod_{j=0}^{e_{\rm int}-1}(aN-2j-1)^{-2}
                         (bN)_{e_{\rm cross}}^{-1}.
\end{equation}
There is no orientation factor for the ordinary internal pairings.
When $a=0$, omit the internal products and exterior stub factor
directly: $m=b$, $c_H=(b!)^{2b}$ and
$P_N^*=(bN)_{b^2}^{-1}$. No $0^0$ convention is needed.

\subsection{Residual simplicity and the leading constants}
After prescribing one assignment, all remaining factors are independent
uniform configurations. Support residual degrees are zero; exterior
vertices, if present, have internal degree $a-1$ and cross degree
$b-1$; all bulk vertices retain $a,b$. Every prescribed edge touches
a saturated support vertex. Hence a residual edge cannot duplicate
a prescribed edge, and full simplicity is exactly residual simplicity.
All leading data and stub assignments give the same residual lists
up to labels, so their residual simplicity probability $s_N^*$ is
the same finite-$N$ number.

For a residual list let $\gamma$ denote its second factorial moment
sum divided by its stub total. If $\delta\ge0$, those totals and
parameters are
\begin{align*}
 L_a&=aN-tD,&\gamma_a&=(a-1)(1-bt/L_a),\\
 L_b&=bN-2bt,&\gamma_b&=(b-1)(1-bt/L_b).
\end{align*}
If $\delta<0$ and $a\ge1$, they are
\begin{align*}
 L_a&=aN-2as,&\gamma_a&=(a-1)(1-as/L_a),\\
 L_b&=bN-sD,&\gamma_b&=(b-1)(1-as/L_b).
\end{align*}
The internal totals are even: $tD=\delta(\delta+1)+2bt$ in the
first case, and $2as$ is even in the second. The cross totals agree
on both sides and need no evenness assumption. At $a=1$ the internal
residual graph is a perfect matching on $N-2s$ vertices and is always
simple. At $a=0$ the internal factors are absent and the cross
residual list is $b$-regular on $N-b$ vertices per side.

For completeness, bounded-degree configuration defect counts give
internal loop and double-edge means $\gamma/2$ and $\gamma^2/4$,
and cross double-edge mean $\gamma_+\gamma_-/2$. At each fixed
factorial-moment order, disjoint defect prescriptions give the product
limits; overlapping distinct prescriptions have edge excess and lose
a power of $N$. Finite Bonferroni bounds for their zero event give
the simplicity limit. The displayed residual parameters tend to
$a-1,b-1$, just as in the original lists. Thus original and residual
product probabilities have the same positive limit
\begin{equation}\label{ap:eq:simplicity}
 s_N,s_N^*\longrightarrow
 \exp\!\left[-\mathbf1_{a\ge1}\frac{a^2-1}{2}
                         -\frac{(b-1)^2}{2}\right],
 \qquad s_N^*/s_N=1+o(1).
\end{equation}
Absent internal factors contribute one, including at $a=0$.
The degree-one internal factors also contribute one. The ratio is
uniform over all leading data because their residual lists coincide
up to labels. This proves eventual simple completability at all
sufficiently large admissible sizes as well. Conditioning produces
exactly the independent uniform simple factors, since their
distinguished-stub multiplicities are constant.

For one leading data pattern its simple-model probability is exactly
$c_HP_N^*s_N^*/s_N$. When $\delta\ge0$, boundary stub powers cancel
to give
\[
 (1+o(1))(a!b!)^{2t}a^{-t\delta}N^{-e_0}.
\]
When $\delta<0$ and $a\ge1$, they instead give
\[
 (1+o(1))(a!b!)^{2s}b^{-s^2}N^{-e_0}.
\]
At $a=0$ the latter coefficient follows by direct omission of the
empty factors in \eqref{ap:eq:prescribed}. Multiplying by
\eqref{ap:eq:data}, for $X$ counting leading flip sets, proves
\begin{equation}\label{ap:eq:first}
 \mathbb E X\sim C_{a,b}N^{-\nu}
\end{equation}
with precisely the two constants in \eqref{eq:ambiguity-exponent}.
The ratio tending to one, rather than a mere simplicity lower bound,
is required for these coefficients.

\subsection{All overlaps and the multiple-state error}
Every possible leading motif is connected and has minimum degree
at least two. For $\delta>0$ its two clique cores are joined by
exterior two-edge paths. For $\delta=0$ it is $K_{2,2a}$. For
$\delta<0$ its $K_{s,s}$ core is connected, with all exterior vertices
attached. At $a=0$ the whole motif is an isolated $K_{b,b}$.
Support vertices have motif degree $D>2$ and exterior vertices degree
two; at $a=0$ every motif vertex is support.

If $H$ is connected with minimum degree at least two, every proper
subgraph $F$ with nonempty vertex set satisfies
\[
 e(F)-v(F)\le e(H)-v(H)-1.
\]
A spanning proper subgraph omits an edge. Otherwise remove $r>0$
vertices. If $E_R$ edges lie among them and $B$ leave them, then
$B>0$ and $2E_R+B\ge2r$, so $E_R+B>r$. At least $r+1$ edges
are removed; further deletion only lowers excess. This includes
retained isolated vertices.

For two distinct leading occurrences include every common vertex
and edge in their intersection $F$. If nonempty it is proper: an
identical full image would have the same motif-degree-$D$ support
and the same uniquely determined boundary data. Distinct data on one
support cannot co-occur because support stubs are saturated. Their
union therefore has excess at least $\nu+1$. Each of the finitely
many typed overlap patterns contributes $O(N^{v-e})$ by its finite
stub count and bounded-cost simplicity conditioning. Incompatible
patterns contribute zero. Disjoint occurrences contribute
$O(N^{-2\nu})=O(N^{-\nu-1})$, since $\nu\ge1$. In particular the
exponent-one case $(1,2)$ requires no stronger disjoint bound.
Consequently
\begin{equation}\label{ap:eq:second}
 \mathbb E[X(X-1)]=O(N^{-\nu-1}),\qquad
 \Prob(X\ge1)=\mathbb E X+O(N^{-\nu-1}),\qquad
 \Prob(X\ge2)=O(N^{-\nu-1}).
\end{equation}
Outside the larger-flip and nonleading-minimum exceptional events,
every alternative is uniquely a leading minimum flip, so $K(G)=1+X$.
Equations \eqref{ap:eq:large}--\eqref{ap:eq:second} prove
\eqref{eq:sharp-ambiguity}, and hence Theorem~\ref{thm:boolean}.
The multiple-state error does not require a rate for the $o(1)$
simplicity ratio in \eqref{ap:eq:simplicity}.

\subsection{Posterior, exact listing, and a charged polynomial cap}\label{sec:enumeration}
For a fixed valid plant every compatible labeled graph has likelihood
\[
 [G_a(N)^2B_b(N)]^{-1},
\]
independent of the partition, and incompatible graphs have likelihood
zero. The three-factor decomposition is unique once the partition is
specified. Here $G_0(N)=1$ and $G_1(N)$ is the admissible matching
count. A uniform hidden balanced unordered plant therefore has an
exactly uniform posterior on its $K(G)$ valid partitions. Relabeling
invariance identifies its ambiguity probability with the fixed-plant
one. Its optimal conditional error is $1-1/K(G)$, attained by any
verified valid partition, and
\[
 0\le R_N-\tfrac12\Prob(K\ge2)
                  \le\tfrac12\Prob(K\ge3)=O(N^{-\nu-1}).
\]
This proves the Bayes constant $C_{a,b}/2$ without a size-bias or
orientation correction. A publicly supplied deterministic plant
would not support this statistical claim.

Theorem~\ref{thm:frame-bounded-moment} applies to the actual two-class
frame for every parameter pair here. At one fixed odd prime it bounds
the scalar kernel moment at $\delta$, so
Corollary~\ref{fk:cor:moment} gives complete sign listing in
unconditional expected nearly quadratic bit work and nearly linear
workspace. The integer equation rejects modular artifacts, also when
the prime divides $b$; every passing vector is automatically balanced.
Select one representative of each global-sign pair. Returning any
listed partition attains the exact conditional Bayes optimum on every
graph. This proves the Las Vegas assertion of
Theorem~\ref{thm:sharp-ambiguity}.

For the deterministic estimator first check that the input is a
simple $D$-regular graph, returning failure otherwise. Use the bounded
binary moment from the same core theorem and encode signs by
$\tau_i=(-1)^{z_i}$.
Every valid assignment solves
\[
 (A+DI)z=(b\bmod2)\one\quad\text{over }\mathbb F_2,
\]
whose homogeneous matrix is $A-\delta I$ modulo two because
$D+\delta=2a$. Solve this affine system by dense binary elimination.
Cap its cardinality at $n^{\nu+1}$; the core moment and Markov give
interruption probability $O(N^{-\nu-1})$. Otherwise enumerate in
Gray order, checking exact integer neighbor counts, and choose any
valid output. On completion the plant is present in the search.
Elimination costs $O(n^3)$, and generation and checks cost $O_{a,b}(n)$
per candidate, giving bounded $O_{a,b}(n^3+n^{\nu+2})$ elementary
work and $O(n^2)$ storage. On interruption use an arbitrary output.
This raises Bayes risk by at most $O(N^{-\nu-1})=o(N^{-\nu})$ and
therefore preserves its leading constant. It completes the proof of
Theorem~\ref{thm:sharp-ambiguity}.

The sharp constants concern independent random cross factors and the
unconditioned simple-factor law. At $a=0$ the leading ambiguity is an
isolated $K_{b,b}$, so connected conditioning can remove it. No
prescribed-cross analogue of these constants, real-corank asymptotic,
growing-parameter or noisy-recovery conclusion is asserted.
